\chapter{考察}
提案したシステムについては，収集対象のシステムの制約を理解した上で，収集から可視化を可能にできたと考える．
また，ピーク検出手法については，最大震度に比例して反応速度が速くなると考えていたが，結果ではそうではなかった．
震源地から本土，人口密度の高い地域の距離にも影響し，対象とした条件により地震数が限られたが，一定の傾向を捉えられたと考える．
感情分類手段の比較では，当初の予想に反して相関が低い結果となった．これは手段ごとに各感情の概念がそれぞれで異なることや，短文コメントに適していない可能性がある．
